\chapter{Serving HTTP}
\label{ch:serving}

\standfirst{A web server that runs one request per core, and a JSON-RPC
protocol on it whose services are files that load themselves.}

Three categories, one on top of the other. \ct{Network} is TCP sockets.
\ct{JSON-RPC-Server} is an HTTP/1.1 server---a listener, a process per
connection, a request parser, files under a document root---which is what
Tomcat was to Kiss. \ct{Rest-Server} is Kiss's protocol on it: \ct{POST /rest}
with a JSON document naming a class and a method, sessions, one database
transaction per request. And \ct{Tonel}, which the last of those leans on,
reads a class from a file into the running image and writes it back.

\section{A first server}

\begin{st}
| server |
server := HttpServer new port: 8080; documentRoot: 'frontend'.
server at: '/hello' put: [:request | HttpResponse text: 'hello, ', request path].
server start
\end{st}

\ct{at:put:} maps a path to a handler---a block taking the request, or any
object answering \ct{handle:} with an \ct{HttpResponse}. The match is exact.
A path no handler claims is looked for under the document root, and is 404
otherwise. \ct{stop} closes the listener; connections in flight finish what
they are doing.

A port of 0 asks the system for one, and \ct{server port} answers it once the
server has started, which is how the tests avoid stepping on each other.

\subsection{Requests and responses}

\begin{center}
\small
\begin{tabular}{@{}ll@{}}
\toprule
\ctp{method} \ctp{path} \ctp{version} & \ct{'GET'}, \ct{'/index.html'} decoded, \ct{'HTTP/1.1'}\\
\ctp{headerAt:} & by name in any case; nil when absent\\
\ctp{body} & a String of bytes, whole\\
\ctp{parameterAt:} \ctp{query} & the query string, decoded\\
\ctp{contentType} \ctp{isJson} \ctp{isMultipart} & the media type without its parameters\\
\ctp{parts} \ctp{formParameters} \ctp{partNamed:} & a multipart body, taken apart into \ct{HttpPart}s\\
\ctp{remoteAddress} & the peer, as text\\
\bottomrule
\end{tabular}
\end{center}

\begin{st}
HttpResponse json: aJsonObject asJsonString.
HttpResponse ok: '<html>...</html>'.
HttpResponse bytes: pngBytes contentType: 'image/png'.
HttpResponse error: 404 reason: 'not here'.
(HttpResponse status: 201) body: '...'; contentType: 'text/plain'; headerAt: 'Location' put: '/x'
\end{st}

\ct{Content-Length} is always sent and the body is always whole, so a client
knows where a response ends and the connection can carry the next request.

\begin{stnote}[What the parser refuses]
The stream is a stranger's, so every reading is bounded. A line over 8,192
bytes, more than a hundred headers, a body over the server's limit (64 MB
unless \ct{maxRequestBytes:} says otherwise) are refused before they are
read, each with its own status; a body with no declared length---chunked---is
refused with 411, because a JSON-RPC client always knows how long its document
is. A handler that raises \ct{HttpError} answers with that status; a handler
that raises anything else answers 500, and the fault goes to standard error,
not to the client.
\end{stnote}

\section{One process per connection}

Every connection the listener accepts is handed to a new \ct{Process}, and
that process runs on whichever worker is idle. That is the whole of the
threading: with the pool running, \(N\) connections are served on \(N\) cores.

What makes that affordable is that a process waiting for its client to send
something costs no worker at all. A socket here is never read by a blocking
call; a read that cannot complete answers \ct{false}, the process \emph{arms}
the socket and waits on a \ct{Semaphore}, and a thread of the VM's own polls
every armed socket and signals the semaphore when the socket is ready. A
hundred idle keep-alive connections are a hundred parked processes. The
database layer, by contrast, parks the worker inside the driver, because a
query has no non-blocking form; a socket does.

\begin{stwarn}[A socket has one reader]
Two processes receiving from one socket would both arm and both wait, and one
signal wakes one of them. Sharing a socket between processes is a program
error, and \ct{Socket} does not try to make it work. Closing a socket a
process is waiting on wakes that process, which then finds the socket closed.
\end{stwarn}

\begin{stwarn}[whileTrue: bodies are inlined]
A temporary declared inside a \ct{whileTrue:} body is one variable for every
iteration, so a process forked in one iteration that captured it reads
whatever the next iteration stored. \ct{HttpServer}'s accept loop hands each
socket to its process as a \emph{method argument} for exactly this reason.
The gate on the socket layer found two processes reading one socket before
it did.
\end{stwarn}

Under \ct{-bootstrap -tests} there is one interpreter and every process is
green; the same code runs and the same tests pass, because the arm-wait-retry
contract is the same either way. Chapter~\ref{ch:parallel} says why the
timing is the only difference.

\section{Sockets by hand}

Nothing above the socket layer is needed to talk to a network. A client is
a \ct{Socket}, a server is a \ct{ServerSocket}, and a \ct{SocketStream}
reads either by the line:

\begin{st}
socket := Socket connectTo: 'example.com' port: 80.
stream := SocketStream on: socket.
stream nextPutAll: 'GET / HTTP/1.0'; cr; cr; flush.
[stream atEnd] whileFalse: [Transcript show: stream nextLine; cr].
socket close
\end{st}

\begin{st}
listener := ServerSocket listenOn: 0.        "0: the system picks; localPort says which"
[[:c | [c send: (c receiveUpTo: 100)] ensure: [c close]] value: listener accept.
 true] whileTrue: []
\end{st}

\ct{connectTo:port:} waits until the connection is made and raises
\ct{NetError}, carrying the system's own words, if the host cannot be reached
or refuses. \ct{receiveUpTo:} answers what has arrived---at least one byte,
at most the number asked for---and an empty String when the peer has closed;
\ct{send:} sends the whole of its argument, however many pieces the kernel
takes it in. Bytes are bytes: nothing translates a line ending, and
\ct{nextLine} drops the CR before an LF because that is what every text
protocol sends.

Every wait is the process's, not the thread's. A socket with nothing to say
parks the process on a Semaphore that the VM's network thread signals when
something arrives, and the worker runs somebody else meanwhile; a hundred
idle connections are a hundred parked processes and no thread at all.
\ct{readTimeout:} on a \ct{SocketStream} bounds every wait for data and
raises \ct{TimedOut} when it runs out, which is how a server stops waiting on
a client that will never speak again. Closing a socket wakes whoever is
waiting on it, so they find it closed rather than sleeping for ever; and one
process reads a socket at a time---two waiting on one connection is a
program error, not a case the class tries to make work.

\section{A REST server}

\begin{st}
RestServer new
    port: 8080;
    backendDirectory: 'backend';
    documentRoot: 'frontend';
    database: 'DSN=shop;UID=app;PWD=secret';
    start
\end{st}

A service is a Tonel file under the back-end directory, a subclass of
\ct{RestService}:

\begin{plain}
backend/services/Adder.class.st

Class { #name : 'Adder', #superclass : 'RestService', #category : 'services' }
Adder >> addNumbers [
    outjson at: 'num3' put: (injson integerAt: 'num1') + (injson integerAt: 'num2')
]
\end{plain}

and a request is a \ct{POST} to \ct{/rest}:

\begin{plain}
$ curl -d '{"_class":"services.Adder","_method":"addNumbers","num1":3,"num2":4}' \
       http://localhost:8080/rest
{"num3":7,"_Success":true,"_ErrorCode":0,"_BootId":"fcec7be8-..."}
\end{plain}

The method named by \ct{_method} is sent, as it is, to a \emph{new} instance
of the class, with four instance variables already set: \ct{injson}, the
request's document; \ct{outjson}, the reply, empty, which the method fills
in; \ct{db}, the request's \ct{DbConnection}; and \ct{request}, everything
else about the request. Kiss hands the same four to a Java method as
parameters. Here they are state, on an instance that belongs to nobody else,
which on this system is the property that matters: nothing is shared between
requests unless a service chooses to share it.

Only a unary method defined in a subclass of \ct{RestService} can be reached
from outside. \ct{printString}, \ct{halt}, the package's own \ct{requireLogin}
and anything with a colon in it answer \emph{No such method}.

\subsection{The wire}

The protocol is Kiss's, unchanged, so that a Kiss front end's \ct{Server.js}
works as it is. The reply is always HTTP 200 and always JSON:

\begin{center}
\small
\begin{tabular}{@{}ll@{}}
\toprule
\ctp{\_class} \ctp{\_method} \ctp{\_uuid} & in: the service, the method, the session\\
\ctp{\_Success} & out: true or false\\
\ctp{\_ErrorCode} & 0; 2 for \emph{log in first}; $-1$ or the service's own\\
\ctp{\_ErrorMessage} & out, when \ct{_Success} is false\\
\ctp{\_BootId} & out: a uuid made at start, so a front end sees a restart\\
\bottomrule
\end{tabular}
\end{center}

Three methods on the empty class are the framework's: \ct{LoginRequired},
\ct{Login} with \ct{username} and \ct{password}, and \ct{Logout}. \ct{Login}
answers \ct{uuid}, which every later request carries as \ct{_uuid}. The
checking is the application's, in \ct{backend/Login.class.st}:

\begin{st}
Login class >> login: username password: password outjson: outjson request: request
    ...answer a RestUserData from request server userCache newUser:..., or nil

Login class >> checkLogin: aUser request: request
    ...answer whether the user is still allowed in; asked again after two minutes
\end{st}

Login is required when the configuration says so, and otherwise whenever
there is a database. \ct{allowWithoutAuthentication:method:} exempts a method
---a public catalogue, a sign-up form.

Inside a service the user is \ct{request userData}: nil when nobody is
logged in, otherwise a \ct{RestUserData} answering \ct{username},
\ct{userId} and \ct{uuid}, and keeping whatever the application puts beside
them:

\begin{st}
Cart >> add
    | user |
    user := request requireLogin.              "raises, code 2, if there is no user"
    user at: #cart put: (user at: #cart ifAbsent: [OrderedCollection new]), ...
\end{st}

That dictionary is the user's own, shared between every request that user
makes---which on this system means between workers---so every \ct{at:} and
\ct{at:put:} takes its Mutex. \ct{request isLoggedIn} asks without raising.
A user is forgotten by \ct{Logout}, or after \ct{userInactiveSeconds} of
silence, and \ct{Login class>>checkLogin:request:} is asked again every two
minutes whether they are still welcome.

\subsection{Errors, uploads, binary replies}

\begin{st}
RestUserError code: 7 signal: 'no such product'     "the user's news; not logged"
RestLogError new signal: 'stock is low'             "a line in the log"
RestServerError new signal: 'the cache is corrupt'  "a fault; logged with everything"
\end{st}

Any other \ct{Error} inside a service is a fault. The client sees the message
and the code, $-1$ unless the service said otherwise; the log sees what the
tier says. The codes are the reply's \ct{_ErrorCode}: 0 for success, 2 when
a login is required (\ct{RestLoginRequiredError}, which \ct{requireLogin}
raises), $-1$ for anything the service did not number, and
\ct{RestUserError code: 7 signal:} for the numbers a front end switches on.
Whatever the tier, the transaction is rolled back.

An upload arrives as a multipart form, files as \ct{_file-0}, \ct{_file-1}
and so on, and the request answers \ct{uploadFileCount},
\ct{uploadFileName:}, \ct{uploadBytes:} and \ct{saveUploadFile:to:}. A
service that wants to send a file back sends \ct{self returnBinary: bytes}:
the reply is then \ct{application/octet-stream}, the JSON, one byte 3, and the
bytes, which is how Kiss's \ct{Server.binaryCall} expects it.

\subsection{One transaction per request}

A connection is taken from a pool before the service runs and put back after.
The method runs inside a transaction that is committed when it returns and
rolled back when it raises. A service that wants the connection gone sooner
sends \ct{request closeConnection: aBoolean}---commit or roll back---and
nothing more is done to it. The pool is a \ct{SharedQueue} of open
connections, which bounds how many requests are inside the database at once
to the pool's size; \ct{lib/Database} asks that a connection belong to one
process at a time, and the queue is what enforces it.

\section{Services are files}

\ct{services.Adder} names \ct{backend/services/Adder.class.st}. The first
request for it reads the file into the image, and every later request
compares the file's time and size with what was loaded and reads it again if
either changed. Edit the file on a running server and the next call runs the
new code, with a whole class as the unit.

The other direction is the Browser. A method accepted or removed in it, on a
class that came from a file, writes the file---Chapter~\ref{ch:keeping} has
the mechanism. The file on disk is the class.

\begin{stnote}
Two things the image's own compiler could not do until this needed them. It
did not read \ct{:=}---every file in \ct{lib/} is spelt that way, and every
one was compiled by the C bootstrap, so nothing had ever compiled that
spelling \emph{inside} the image---and a compile error with nobody to tell
opened 1983's \ct{SyntaxError} notifier and suspended the process, which in
a headless server is a hang. Both are fixed: \ct{Scanner} reads \ct{:=}, and
a compile error in a file is a \ct{TonelError} naming the file, the method
and the character.
\end{stnote}

\section{Running a server image}

Nothing above needs \ct{-serve}: a \ct{RestServer} started from a workspace
in \ct{-run} serves requests green, one interpreter switching between them.
For a server that uses every core:

\begin{plain}
./st80 -bootstrap -profile profiles/st2026.profile -startup 'RestServer serve' -o server.im
./st80 -serve server.im -workers 8 server.json
\end{plain}

\ct{-serve} loads the image, starts the worker pool, and resumes the image's
startup process on the first worker while the rest join the scheduler with
nothing to run and take processes as they appear. \ct{RestServer class>>serve}
reads its configuration from the file named by the first argument---
\ct{server.json} by default:

\begin{plain}
{ "port": 8080, "backendDirectory": "backend", "documentRoot": "frontend",
  "database": "DSN=shop;UID=app;PWD=secret", "requireAuthentication": true,
  "userInactiveSeconds": 3600, "connectionPoolSize": 12,
  "allowWithoutAuthentication": [ "services.Catalog:list" ] }
\end{plain}

Every name in the file is also in \ct{request environment}, so a service can
read its own settings from the same place. The server stays up until
\ct{SIGINT} or \ct{SIGTERM}, then stops and quits with exit code 0. The log is
standard error, one line per event.

\begin{stnote}
The startup is baked in at bootstrap and not given to \ct{-serve}, because
the bootstrap's compiler resolves globals through tables a loaded image does
not have. A desktop image run under \ct{-serve} executes its display loop
headless on one worker and serves nothing.
\end{stnote}

\section{Where the rest of it is}

\ct{doc/NETWORK.md} is the socket layer and the run mode: the one decision
everything follows from, the arm-wait-retry contract, what the VM does, and
the two things found the hard way. \ct{doc/REST-SERVER.md} is the protocol,
where Kiss stops and this begins, and what is shared and what guards it.
\ct{lib/Rest-Server/PROVENANCE.md} records what of Kiss crossed, what changed,
and the three Kiss faults not repeated. \ct{doc/CONCURRENCY.md} records what
the gates found in the virtual machine itself while this was built---a list
unlinked in the wrong order, a slot every worker wrote, an enumeration that
walked headers being refilled---each older than the server and each waiting
for thirty-one workers to show it.

The tests are in \ct{lib/Network-Tests}, \ct{lib/JSON-RPC-Server-Tests} and
\ct{lib/Rest-Server-Tests}, all over the loopback interface on a port the
system picks; \ct{lib/Rest-Server-Live-Tests} needs a database and lives in
the \ct{database-live} profile. The gates are
\ct{tests/unit/test_parallel_net.c} and
\ct{tests/unit/test_parallel_rest.c}: thirty-one workers, sixty-two native
clients, every sum checked, every worker seen.
